\begin{table*}[!t]
\centering\small
\begin{tabular}{@{}p{.19\textwidth}p{.69\textwidth}r@{}}
\toprule
Questionnaire & Item & $\pi_i$ \\
\midrule
HEXACO & I think that most people like some aspects of my personality. & 18.38 \\
BFI-2 & Is curious about many different things. & 18.00 \\
FFMQ & I don’t pay attention to what I’m doing because I’m daydreaming, worrying, or otherwise distracted. & 16.06 \\
Schwartz Values & I seek out pleasure in life & 15.82 \\
BFI-2 & Is persistent, works until the task is finished. & 15.69 \\
MLQ & I am seeking a purpose or mission for my life. & 15.22 \\
IRI & When I am reading an interesting story or novel, I imagine how I would feel if the events in the story were happening to me. & 13.90 \\
Internal Representation & I can close my eyes and easily picture a scene that I have experienced & 13.53 \\
HEXACO & When someone I know well is unhappy, I can almost feel that person's pain myself. & 13.33 \\
DERS & I care about what I am feeling & 13.31 \\
BFI-2 & Is compassionate, has a soft heart. & 13.22 \\
Schwartz Values & I value my creativity and individuality & 13.03 \\
Internal Representation & When I read, I tend to hear a voice in my "mind’s ear" & 12.92 \\
ECR-R & It helps to turn to my romantic partner in times of need. & 12.82 \\
BFI-2 & Feels little sympathy for others. & 12.75 \\
BIS/BAS & How I dress is important to me. & 12.63 \\
VISQ-R & I evaluate my behavior using my inner speech. For example I say to myself, “that was good” or “that was stupid” & 12.38 \\
Self-Monitoring & In conversations, I am sensitive to even the slightest change in the facial expression of the person I’m conversing with. & 12.30 \\
MCQ-30 & I am aware of the way my mind works when I am thinking through a problem & 12.27 \\
EPQ-R & Are you sometimes bubbling over with energy and sometimes very sluggish? & 12.06 \\
\bottomrule
\end{tabular}
\caption{The 20 highest neutral-to-human-simulation variance ratios among
1,312 eligible items. Selection uses the untransformed ratio, with at least
five valid responses in each condition and positive human-simulation variance.
The ratio measures relative disagreement, not the direction of an item's
association with PC1.}
\label{tab:ratio_item_examples}
\end{table*}
